% TOP_PROBLEM: {"id": 9400155, "title": "Are there infinitely many palindromic primes to every base", "category_id": 1, "category": {"id": 1, "name": "number_theory", "display_name": "Number Theory", "description": "Properties of integers, prime numbers, Diophantine equations.", "slug": "number-theory", "order_index": 1, "created_at": "2026-07-31T15:26:25.670Z"}, "status": "partially_solved"}
% FIRST_POSTED: null
% ENTRY_KIND: statement_only
% Imported from: batch12.tex; UnsolvedMath ID 9400155
\documentclass[11pt]{article}
\usepackage[a4paper,margin=27mm]{geometry}
\usepackage{amsmath,amssymb,mathtools}
\usepackage{fontspec,unicode-math}
\setmainfont{Latin Modern Roman}
\setmathfont{Latin Modern Math}
\usepackage{xurl}
\usepackage[hidelinks]{hyperref}
\newcommand{\sref}[2]{\hyperlink{source:#1:#2}{[S#2]}}
\newcommand{\cataloguescope}[1]{\paragraph{Mathematical question.}#1}
\hypersetup{pdftitle={UnsolvedMath Problem 9400155}}
\begin{document}
% UnsolvedMath ID 9400155; source catalogue code AMR-093-0155
\setcounter{section}{9400154}
\section{Are there infinitely many palindromic primes to every base}\label{9400155}
{\small\sffamily UnsolvedMath ID 9400155 · Number Theory}
\subsection{Definitions and mathematical statement}

\paragraph{Shared notation.}$\mathbb N=\{1,2,\ldots\}$, and $\mathbb Z,\mathbb Q,\mathbb R,\mathbb C$ denote the integers, rationals, reals and complex numbers. Any other conventions are specified locally.


For an integer base $b\ge2$, write the unique base-$b$ expansion of $N\ge1$ as $N=\sum_{j=0}^m a_jb^j$, where $m\ge0$, $0\le a_j<b$ and $a_m\ne0$. It is palindromic in base $b$ if $a_j=a_{m-j}$ for every $0\le j\le m$. A base-$b$ palindromic prime is a prime with this digit symmetry.
\paragraph{Statement recorded in the supplied source.}
For every integer $b\ge2$, are there infinitely many primes palindromic in base $b$? In particular, is this true for base $10$? \sref{9400155}{2}
\subsection{Short English statement}
Does each integer base have infinitely many primes whose digits read the same in both directions?
\subsection{Sources}
\begin{description}
\item[\hypertarget{source:9400155:1}{S1}] ULAM.AI and the UnsolvedMath Contributors, UnsolvedMath: A Curated Collection of Open Mathematics Problems, record 9400155 (AMR-093-0155). CC BY 4.0. \url{https://huggingface.co/datasets/ulamai/UnsolvedMath}.
\item[\hypertarget{source:9400155:2}{S2}] Wikipedia contributors, \emph{List of unsolved problems in mathematics}, revision 1366636767 (2026) Number theory, Are there infinitely many palindromic primes to every base. \url{https://en.wikipedia.org/w/index.php?title=List_of_unsolved_problems_in_mathematics&oldid=1366636767#Number_theory}.
\item[\hypertarget{source:9400155:3}{S3}] Phakhinkon Phunphayap and Prapanpong Pongsriiam, \emph{Reciprocal Sum of Palindromes}, arXiv:1803.00161v2 (2018) Abstract and Section 1, p. 1. \url{https://arxiv.org/pdf/1803.00161}.
\end{description}
\label{end:9400155}
\end{document}
